\casechapter{Inequalities as Maps}{Optimal transport}{Geometrization}{ch:transport}
\begin{refsection}[chapters/10-transport.bib]

% Notation for this chapter: relative entropy, relative Fisher information,
% the quadratic transport distance, push-forward, and a class of regular measures.
\newcommand{\transportH}[2]{H(#1\,|\,#2)}
\newcommand{\transportI}[2]{I(#1\,|\,#2)}
\newcommand{\transportW}{W_2}
\newcommand{\transportpf}{_{\#}}
\newcommand{\transportReg}{\mathcal{R}}

\begin{chaptersummary}
The Brunn--Minkowski, isoperimetric, and Pr\'ekopa--Leindler inequalities, and the Gaussian logarithmic Sobolev and transport inequalities, were proved by unrelated devices: cutting boxes, symmetrizing, slicing by level sets, passing through the discrete cube. The move is to carry one measure onto the other by a map, so that the inequality becomes a pointwise inequality for the Jacobian, integrated once. Taken further, the maps make the space of probability measures a geometry in which these inequalities are elementary facts about convex functions. What remains hard is constructing the maps.
\end{chaptersummary}

\section{Before: one device per inequality}\label{sec:transport:before}

The inequalities of this chapter compare quantities attached to sets, functions, or measures. For most of the twentieth century each was proved by its own device.

\subsection{Brunn--Minkowski: cutting boxes}\label{sec:transport:bmbefore}

For $A,B\subseteq\R^n$ let $A+B=\{a+b:a\in A,\ b\in B\}$ be their \term{Minkowski sum}, and let $\abs{A}$ be Lebesgue measure.

\begin{theorem}[Brunn--Minkowski]\label{thm:transport:bm}
For all nonempty compact sets $A,B\subseteq\R^n$,
\begin{equation}\label{eq:transport:bm}
\abs{A+B}^{1/n}\ \ge\ \abs{A}^{1/n}+\abs{B}^{1/n}.
\end{equation}
\end{theorem}

Equality holds for homothetic convex bodies.\index{Brunn--Minkowski inequality} Brunn found the inequality for convex bodies in three dimensions in his 1887 thesis, Minkowski proved it for convex bodies in all dimensions, and Lusternik extended it to measurable sets in 1935.\sidenote{Minkowski also found an error in Brunn's argument, which Brunn corrected. \textcite{Gardner2002} tells the whole story.} The shortest classical proof is due to \textcite{HadwigerOhmann1956}.

\begin{aside}[The Hadwiger--Ohmann proof]
\emph{Boxes.} If $A$ and $B$ are boxes with side lengths $a_i$ and $b_i$, then $A+B$ has side lengths $a_i+b_i$, and the arithmetic--geometric mean inequality gives
\[
\prod_{i=1}^n\Bigl(\frac{a_i}{a_i+b_i}\Bigr)^{1/n}+\prod_{i=1}^n\Bigl(\frac{b_i}{a_i+b_i}\Bigr)^{1/n}\le\frac1n\sum_{i=1}^n\frac{a_i}{a_i+b_i}+\frac1n\sum_{i=1}^n\frac{b_i}{a_i+b_i}=1.
\]
\emph{Unions of boxes.} Induct on the number of boxes. Translate $A$ so that the hyperplane $\{x_n=0\}$ separates two of its boxes, and translate $B$ along the $x_n$-axis until its two sides $B^\pm$ carry the same fractions of its volume as the two sides $A^\pm$ of $A$. Each pair $(A^\pm,B^\pm)$ has fewer boxes, and $A^++B^+$ and $A^-+B^-$ lie on opposite sides of the hyperplane, so by induction $\abs{A+B}\ge\sum_\pm\bigl(\abs{A^\pm}^{1/n}+\abs{B^\pm}^{1/n}\bigr)^n=\bigl(\abs{A}^{1/n}+\abs{B}^{1/n}\bigr)^n$.
\emph{Compact sets.} Approximate from outside by unions of dyadic cubes.
\end{aside}

Everything here is about boxes, and nothing suggests how to prove an inequality of the same shape on a sphere, on a curved manifold, or for another measure.

\subsection{Isoperimetry: symmetrize, and hope that a minimizer exists}\label{sec:transport:isobefore}

The \term{isoperimetric inequality} says that among sets of given volume the ball has the least surface area.\index{isoperimetric inequality} In 1838 Steiner argued by \term{Steiner symmetrization}:\index{Steiner symmetrization} sliding each chord in a fixed direction until its midpoint lies on a fixed line orthogonal to the chords preserves area, does not increase perimeter, and strictly improves, for a suitable direction, every convex figure other than a disc. Hence the disc is optimal \emph{if} an optimal figure exists. That is not a formality: Weierstrass showed in 1870 that a functional bounded below need not attain its infimum.\sidenote{Weierstrass's example was aimed at Riemann's use of the Dirichlet principle, which assumed that a certain energy is minimized by some function.} He treated the isoperimetric problem in his lectures on the calculus of variations, and Steiner's argument was completed by compactness arguments such as Blaschke's selection theorem.\sidenote{Blaschke proved the selection theorem in his 1916 book \emph{Kreis und Kugel} (Circle and Sphere).} The price of symmetrization is an existence theorem.

\subsection{Pr\'ekopa--Leindler: slice, then induct}\label{sec:transport:plbefore}

Pr\'ekopa (1971) and Leindler (1972) proved Brunn--Minkowski for functions.\index{Pr\'ekopa--Leindler inequality}

\begin{theorem}[Pr\'ekopa--Leindler]\label{thm:transport:pl}
Let $0<t<1$, and let $f,g,h$ be nonnegative measurable functions on $\R^n$, with $f$ and $g$ integrable, such that
\begin{equation}\label{eq:transport:plhyp}
h\bigl((1-t)x+ty\bigr)\ \ge\ f(x)^{1-t}g(y)^{t}\qquad\text{for all } x,y\in\R^n.
\end{equation}
Then $\int h\ge\bigl(\int f\bigr)^{1-t}\bigl(\int g\bigr)^{t}$.
\end{theorem}

Indicators of compact sets give $\abs{(1-t)A+tB}\ge\abs{A}^{1-t}\abs{B}^{t}$, which implies \eqref{eq:transport:bm} by rescaling. The standard proof normalizes $\sup f=\sup g=1$, applies $\abs{X+Y}\ge\abs{X}+\abs{Y}$ on the line to the level sets $\{h\ge s\}\supseteq(1-t)\{f\ge s\}+t\{g\ge s\}$, integrates over $s$, and inducts on dimension through slices. A second one-dimensional argument, used by Henstock and Macbeath (1953) for Brunn--Minkowski, parametrizes by mass instead of by height:\sidenote{See \textcite{Gardner2002}. In one dimension the move of this chapter is old; what was missing was a way to use it in all dimensions at once, and an explanation of why it works.} the one-dimensional case of the move of this chapter.

\subsection{Gaussian inequalities: the cube, the semigroup, and tensorization}\label{sec:transport:gaussbefore}

Let $\gamma$ be the standard Gaussian measure on $\R^n$. For probability measures $\nu$ and $\pi$, the \term{relative entropy} is $\transportH{\nu}{\pi}=\int\log\frac{\dd\nu}{\dd\pi}\dd\nu$ ($+\infty$ unless $\nu$ is absolutely continuous with respect to $\pi$), and, when $f=\dd\nu/\dd\pi$ is positive and differentiable, the relative \term{Fisher information} is $\transportI{\nu}{\pi}=\int\abs{\nabla\log f}^2\dd\nu$. These normalizations hold throughout.

\begin{theorem}[\textcite{Gross1975}]\label{thm:transport:gross}
For every probability measure $\nu$ on $\R^n$ whose density with respect to $\gamma$ is positive and continuously differentiable,
\begin{equation}\label{eq:transport:lsi}
\transportH{\nu}{\gamma}\le\tfrac12\,\transportI{\nu}{\gamma}.
\end{equation}
\end{theorem}

With $u=\sqrt{\dd\nu/\dd\gamma}$ it reads $\int u^2\log u^2\dd\gamma\le2\int\abs{\nabla u}^2\dd\gamma$, Gross's form of the \term{logarithmic Sobolev inequality}.\index{logarithmic Sobolev inequality}\sidenote{Gross worked in constructive quantum field theory, and he showed that the inequality is equivalent to the hypercontractivity of the Ornstein--Uhlenbeck semigroup; compare the role of hypercontractivity on the cube in \cref{ch:fourier}.} His proof went through the discrete cube.

\begin{aside}[Gross's route]
\emph{Two points.} For the uniform measure on $\{-1,1\}$, prove by calculus that $\operatorname{Ent}(u^2)\le\frac12(u(1)-u(-1))^2$, where $\operatorname{Ent}(u^2)=\E[u^2\log u^2]-\E[u^2]\log\E[u^2]$.
\emph{The cube.} Entropy is subadditive on products, which gives the inequality on $\{-1,1\}^m$ with a sum over coordinates on the right.
\emph{The central limit theorem.} Apply it to $u(x)=\phi((x_1+\dots+x_m)/\sqrt m)$ and let $m\to\infty$. The result is \eqref{eq:transport:lsi} in dimension one; tensorization gives all dimensions.
\end{aside}

\textcite{BakryEmery1985} found a route without cubes. Call a differentiable $V$ on $\R^n$ \term{$\rho$-convex}, for $\rho>0$, if $V(y)\ge V(x)+\nabla V(x)\cdot(y-x)+\frac{\rho}{2}\abs{y-x}^2$ for all $x,y$. For $\pi=Z^{-1}e^{-V(x)}\dd x$ with $V$ twice differentiable and $\rho$-convex, they differentiated the entropy along the semigroup with generator $\Delta-\nabla V\cdot\nabla$ and proved
\begin{equation}\label{eq:transport:be}
\transportH{\nu}{\pi}\le\frac{1}{2\rho}\,\transportI{\nu}{\pi}
\end{equation}
for every $\nu$ with a regular positive density; the Gaussian is $V(x)=\abs{x}^2/2$, $\rho=1$. The proof is short once the semigroup is under control, and hard to find without it. A third inequality uses the distance $\transportW$: $\transportW(\mu,\nu)^2$ is the least value of $\E\abs{X-Y}^2$ over random vectors $X\sim\mu$ and $Y\sim\nu$ (\cref{def:transport:coupling}).

\begin{theorem}[\textcite{Talagrand1996}]\label{thm:transport:talagrand}
For every probability measure $\nu$ on $\R^n$, $\transportW(\nu,\gamma)^2\le2\,\transportH{\nu}{\gamma}$.
\end{theorem}

Talagrand proved this \term{transport inequality}\index{transport inequality} in dimension one with the monotone map carrying $\gamma$ to $\nu$, a transport proof before the name, and passed to $\R^n$ by tensorization; Marton had found in 1986 that such inequalities imply concentration of measure. Three inequalities, three proofs, each tied to a product structure or an analytic machine. Yet the two Gaussian inequalities are sharp for the same measures, the translates of $\gamma$ (Exercise~\ref{ex:transport:gauss}), and one is led to ask whether, for every probability measure on $\R^n$, \eqref{eq:transport:be} implies Talagrand's inequality in the form $\transportW^2\le\frac2\rho H$.

\section{The move: carry one measure onto the other}\label{sec:transport:move}

Every inequality above compares two objects of the same kind. The move is to compare them point by point, along a map that carries one onto the other: after the substitution $y=T(x)$ the inequality becomes a pointwise inequality for the derivative of $T$, integrated once. Grothendieck's advice to stoop low applies: the tools are the change of variables formula, the arithmetic--geometric mean inequality, and $\log x\le x-1$. The second half of the move came later: the optimal maps make the space of probability measures a geometry, with straight lines, convexity, and gradient flows, in which the inequalities become facts about convex functions.

\subsection{Maps, couplings, and the quadratic distance}\label{sec:transport:couplings}

\begin{definition}[Transport maps, couplings, and $\transportW$]\label{def:transport:coupling}
Let $\mu$ and $\nu$ be Borel probability measures on $\R^n$. A Borel map $T\colon\R^n\to\R^n$ \term{transports} $\mu$ to $\nu$ if $T\transportpf\mu=\nu$, that is, $\mu(T^{-1}(E))=\nu(E)$ for every Borel set $E$. A \term{coupling} of $\mu$ and $\nu$ is a probability measure on $\R^n\times\R^n$ with marginals $\mu$ and $\nu$. On the set $\cP_2(\R^n)$ of probability measures with finite second moments, the \term{quadratic Wasserstein distance} is
\[
\transportW(\mu,\nu)=\Bigl(\inf_{\pi}\int\abs{x-y}^2\dd\pi(x,y)\Bigr)^{1/2},
\]
the infimum being over all couplings $\pi$ of $\mu$ and $\nu$.
\end{definition}

The infimum is attained, and $\transportW$ is a metric \parencite{Villani2003}; the cost $\abs{x-y}$, without square or root, gives the distance $W_1\le\transportW$.\index{optimal transport}\index{Wasserstein distance} A map $T$ transporting $\mu$ to $\nu$ gives the coupling $(\mathrm{id},T)\transportpf\mu$, so
\begin{equation}\label{eq:transport:mapbound}
\transportW(\mu,\nu)^2\le\int\abs{T(x)-x}^2\dd\mu(x).
\end{equation}
\textcite{Monge1781} asked for the cheapest way to move earth from excavations to embankments, and he looked for maps.\sidenote{The title of Monge's memoir names the problem: \emph{d\'eblais} are excavated earth, \emph{remblais} are embankments.} \textcite{Kantorovich1942} replaced maps by couplings, which makes the problem linear.\sidenote{Kantorovich shared the 1975 Nobel Memorial Prize in Economic Sciences with Tjalling Koopmans. The name ``Wasserstein distance'' was coined by Dobrushin in 1970, after a 1969 paper of L.~N.~Vaserstein.}

\begin{theorem}[Kantorovich duality, discrete case]\label{thm:transport:kantorovich}
Let $\mu=\sum_{i=1}^k a_i\delta_{x_i}$ and $\nu=\sum_{j=1}^l b_j\delta_{y_j}$ be probability measures and let $c_{ij}$ be real numbers. The minimum of $\sum_{i,j}c_{ij}\pi_{ij}$ over nonnegative arrays $(\pi_{ij})$ with row sums $a_i$ and column sums $b_j$ equals the maximum of $\sum_ia_i\varphi_i+\sum_jb_j\psi_j$ over real vectors $\varphi,\psi$ with $\varphi_i+\psi_j\le c_{ij}$ for all $i,j$, and both extrema are attained.
\end{theorem}

\begin{proof}
The minimum is a feasible linear program in standard form ($\pi_{ij}=a_ib_j$ is admissible) with a compact feasible set, so it is attained. The maximum is its dual, with one free variable for each equality constraint and one inequality for each variable $\pi_{ij}$, and strong duality of linear programming (\cref{ch:lpduality}) says that it is attained and equal to the minimum.
\end{proof}

\index{Kantorovich duality}For the quadratic cost, with $\Phi_i=\frac12(\abs{x_i}^2-\varphi_i)$ and $\Psi_j=\frac12(\abs{y_j}^2-\psi_j)$, the constraint reads $\Phi_i+\Psi_j\ge x_i\cdot y_j$, and for optimal $\pi$ and $(\varphi,\psi)$ complementary slackness gives equality where $\pi_{ij}>0$; so an optimal plan lives on the graph of the subdifferential of the convex function $\max_j(x\cdot y_j-\Psi_j)$, a discrete shadow of Brenier's theorem below.

\subsection{One dimension: the monotone rearrangement}\label{sec:transport:oned}

For a distribution function $F$ let $F^{-1}(u)=\inf\{x:F(x)\ge u\}$, for $0<u<1$, be its quantile function.

\begin{theorem}[Optimal transport on the line]\label{thm:transport:oned}
Let $\mu$ and $\nu$ be probability measures on $\R$ with finite second moments and distribution functions $F$ and $G$, and let $U$ be uniform on $(0,1)$. The law $\pi^\star$ of $\bigl(F^{-1}(U),G^{-1}(U)\bigr)$ is an optimal coupling for the quadratic cost, so
\begin{equation}\label{eq:transport:w2line}
\transportW(\mu,\nu)^2=\int_0^1\bigl(F^{-1}(u)-G^{-1}(u)\bigr)^2\dd u.
\end{equation}
If $F$ is continuous, then the nondecreasing map $T=G^{-1}\circ F$, defined where $0<F<1$, transports $\mu$ to $\nu$, and $\transportW(\mu,\nu)^2=\int(T(x)-x)^2\dd\mu(x)$.
\end{theorem}

\begin{proof}
The plan is to show that $\pi^\star$ maximizes $\int xy\dd\pi$ over couplings $\pi$; this is the same as minimizing the cost, because $\int(x-y)^2\dd\pi=\int x^2\dd\mu+\int y^2\dd\nu-2\int xy\dd\pi$.

\emph{The coupling.} For $0<u<1$, right-continuity of $F$ gives $F^{-1}(u)\le x$ if and only if $u\le F(x)$, so $\Pr[F^{-1}(U)\le x]=F(x)$; likewise for $G$.

\emph{Hoeffding's identity.} Let $(X,Y)$ have law $\pi$, let $J(s,t)=\Pr[X\le s,\,Y\le t]$, and let $(X',Y')$ be an independent copy of $(X,Y)$. For real $a$ and $b$, $a-b=\int\bigl(\ind\{b\le s\}-\ind\{a\le s\}\bigr)\dd s$. Multiply this identity for $(X,X')$ by the one for $(Y,Y')$, take expectations, and exchange the order of integration, which is legitimate because the absolute value of the integrand integrates to $\abs{X-X'}\abs{Y-Y'}$, a variable with finite mean. The result is
\[
2\Cov(X,Y)=\E\bigl[(X-X')(Y-Y')\bigr]=2\iint\bigl(J(s,t)-F(s)G(t)\bigr)\dd s\dd t.
\]

\emph{The comparison.} Always $J(s,t)\le\min\{F(s),G(t)\}$, and $\pi^\star$ attains equality, since $\Pr[F^{-1}(U)\le s,\,G^{-1}(U)\le t]=\Pr[U\le\min\{F(s),G(t)\}]$. The means are fixed by the marginals, so $\int xy\dd\pi\le\int xy\dd\pi^\star$, which proves optimality and \eqref{eq:transport:w2line}.

\emph{The map.} If $F$ is continuous and $X\sim\mu$, then $U=F(X)$ is uniform on $(0,1)$, and $F^{-1}(F(X))=X$ almost surely, because $F^{-1}(F(x))<x$ only if $F$ is constant on an interval ending at $x$, and such intervals carry no mass. So $(X,T(X))=\bigl(F^{-1}(U),G^{-1}(U)\bigr)$ almost surely.
\end{proof}

\begin{figure}[t]
\centering
\begin{tikzpicture}[x=1.9cm,y=1cm,
  declare function={gs(\m,\s,\t)=exp(-((\t-\m)^2)/(2*\s*\s))/(\s*sqrt(2*pi));
    ftop(\t)=0.45*gs(1.3,0.42,\t)+0.55*gs(3.9,0.55,\t);
    gbot(\t)=gs(2.7,0.85,\t);},
  dens/.style={line width=0.8pt,seadeep},
  link/.style={sea,line width=0.45pt,-{Stealth[length=1.5mm]}}]
\fill[seamist] (1.241,1.9) -- plot[domain=1.241:1.481,samples=12] (\x,{1.9+1.8*ftop(\x)}) -- (1.481,1.9) -- cycle;
\fill[seamist] (1.985,0) -- plot[domain=1.985:2.254,samples=12] (\x,{-1.8*gbot(\x)}) -- (2.254,0) -- cycle;
\draw[muted,line width=0.4pt] (0.2,1.9) -- (5.4,1.9);
\draw[muted,line width=0.4pt] (0.2,0) -- (5.4,0);
\draw[dens] plot[domain=0.2:5.4,samples=150] (\x,{1.9+1.8*ftop(\x)});
\draw[dens] plot[domain=0.2:5.4,samples=150] (\x,{-1.8*gbot(\x)});
\foreach \a/\b in {0.979/1.611,1.241/1.985,1.481/2.254,1.812/2.485,3.166/2.7,3.567/2.915,3.837/3.146,4.092/3.415,4.4/3.789}
  \draw[link] (\a,1.86) -- (\b,0.04);
\node[anchor=west] at (5.4,2.15) {$\mu$};
\node[anchor=west] at (5.4,-0.3) {$\nu$};
\node[sea] at (4.45,0.95) {$T$};
\end{tikzpicture}
\caption{The monotone rearrangement $T=G^{-1}\circ F$ carries each decile of $\mu$ to the same decile of $\nu$. The shaded slabs carry equal mass, so $T$ stretches where $\nu$ is thinner than $\mu$.}
\label{fig:transport:monotone}
\end{figure}

The map $T$ is nondecreasing, hence the derivative of a convex function: in one dimension the optimal map is the gradient of a convex function (\cref{fig:transport:monotone}). The same coupling is optimal for every convex cost $h(x-y)$ (Exercise~\ref{ex:transport:convexcost}). This \term{monotone rearrangement}\index{monotone rearrangement} is the map behind the parametrization of Henstock and Macbeath and behind Talagrand's proof.

\subsection{Higher dimensions: Knothe's map and Brenier's map}\label{sec:transport:ndmaps}

\textcite{Knothe1957} built a map in $\R^n$ by applying the monotone rearrangement one coordinate at a time, to the conditional law of each coordinate given the previous ones.\sidenote{Rosenblatt had used the same triangular construction in statistics in 1952; the map is also called the Knothe--Rosenblatt rearrangement.} The $k$th component of this \term{Knothe map}\index{Knothe map} depends only on $x_1,\dots,x_k$ and increases in $x_k$, so its Jacobian matrix is triangular with positive diagonal. The map is elementary, but it depends on the order of the coordinates, and it is not optimal (Exercise~\ref{ex:transport:knothe}).

\begin{theorem}[\textcite{Brenier1991}; \textcite{McCann1995}]\label{thm:transport:brenier}
Let $\mu$ and $\nu$ be probability measures on $\R^n$, with $\mu$ absolutely continuous. There is a convex function $\varphi\colon\R^n\to\R\cup\{+\infty\}$ whose gradient, which exists $\mu$-almost everywhere, transports $\mu$ to $\nu$, and the map $\nabla\varphi$ is uniquely determined $\mu$-almost everywhere. If $\mu$ and $\nu$ have finite second moments, then $(\mathrm{id},\nabla\varphi)\transportpf\mu$ is the unique optimal coupling, so $\transportW(\mu,\nu)^2=\int\abs{\nabla\varphi(x)-x}^2\dd\mu(x)$.
\end{theorem}

We do not prove \term{Brenier's theorem}\index{Brenier's theorem} \parencite[see][]{Villani2003}; the discrete shadow above shows where $\varphi$ comes from. For $\mu=f\dd x$ and $\nu=g\dd x$ it gives the Monge--Amp\`ere equation $f(x)=g(\nabla\varphi(x))\det D^2\varphi(x)$, valid $\mu$-almost everywhere with the Hessian that a convex function has almost everywhere \parencite{McCann1997}.\sidenote{By Aleksandrov's theorem, a convex function has a second-order Taylor expansion at almost every point.} This Jacobian is symmetric and positive semidefinite.

\begin{definition}[Displacement interpolation and convexity]\label{def:transport:displacement}
Let $\mu_0,\mu_1\in\cP_2(\R^n)$, with $\mu_0$ absolutely continuous, and let $T=\nabla\varphi$ be the map of \cref{thm:transport:brenier} from $\mu_0$ to $\mu_1$. The \term{displacement interpolation} is the curve $\mu_t=\bigl((1-t)\mathrm{id}+tT\bigr)\transportpf\mu_0$, $0\le t\le1$. A functional $\cF$ on $\cP_2(\R^n)$ is \term{displacement $\rho$-convex} if, for all such pairs and all $t\in[0,1]$,
\[
\cF(\mu_t)\le(1-t)\cF(\mu_0)+t\cF(\mu_1)-\tfrac{\rho}{2}\,t(1-t)\,\transportW(\mu_0,\mu_1)^2.
\]
\end{definition}

\index{displacement interpolation}Particles move at constant speed along segments, and $\transportW(\mu_s,\mu_t)=\abs{t-s}\,\transportW(\mu_0,\mu_1)$ \parencite{Villani2003}: the curve is a geodesic. It is not the mixture $(1-t)\mu_0+t\mu_1$, which keeps particles in place and changes their weights; halfway between two distant Gaussians on the line the mixture is bimodal, while the displacement interpolant is the Gaussian with averaged mean and standard deviation. \textcite{McCann1997} introduced displacement interpolation for interacting gases, whose energy is not convex along mixtures, and \textcite{Otto2001}, after \textcite{JordanKinderlehrerOtto1998}, made $\cP_2(\R^n)$ a formal Riemannian manifold with these geodesics.

What had to be noticed? The arithmetic--geometric mean step of Hadwiger and Ohmann is the Jacobian inequality for a map between boxes, and the parametrizations of Henstock--Macbeath and Talagrand are the monotone map. The insight was that good maps exist in every dimension, so no induction is needed, and that their straight lines are geodesics along which entropy is convex.

\section{After: inequalities from Jacobians}\label{sec:transport:after1}

We now prove the Pr\'ekopa--Leindler, Brunn--Minkowski, and isoperimetric inequalities in $\R^n$, and the Gaussian and log-concave inequalities on the line; three lemmas about maps carry the difficulty.

\subsection{Pr\'ekopa--Leindler in all dimensions at once}\label{sec:transport:plafter}

Write $x_{<k}=(x_1,\dots,x_{k-1})$ and $x_{\le k}=(x_1,\dots,x_k)$.

\begin{lemma}[Change of variables for triangular maps]\label{lem:transport:triangular}
Let $Q=\prod_{k=1}^n(a_k,b_k)$ be an open box and $z\colon Q\to\R^n$ a Borel map such that each component $z_k$ depends only on $x_{\le k}$ and, for each fixed $x_{<k}$, the function $s\mapsto z_k(x_{<k},s)$ is continuously differentiable with positive derivative on $(a_k,b_k)$. Then every Borel function $h\colon\R^n\to[0,\infty]$ satisfies
\[
\int_Q h(z(x))\prod_{k=1}^n\partial_kz_k(x)\dd x\ \le\ \int_{\R^n}h(y)\dd y.
\]
\end{lemma}

\begin{proof}
We induct on $n$; the functions $\partial_kz_k$ are Borel, as limits of difference quotients. For $n=1$ the substitution $y=z(x)$ gives $\int_{a_1}^{b_1}h(z(x))z'(x)\dd x=\int_{z((a_1,b_1))}h\le\int_{\R}h$. For $n>1$ write $x=(x',x_n)$ and $z=(z',z_n)$, where $z'$ depends only on $x'$, and let $\bar h(y')=\int_{\R}h(y',y_n)\dd y_n$. For fixed $x'$, the case $n=1$ applied to $s\mapsto z_n(x',s)$ gives
\[
\int_{a_n}^{b_n}h\bigl(z'(x'),z_n(x',s)\bigr)\,\partial_nz_n(x',s)\dd s\ \le\ \bar h(z'(x')).
\]
Multiply by $\prod_{k<n}\partial_kz_k(x')$, integrate over $x'$ (Tonelli's theorem), and apply the induction hypothesis to $z'$ and $\bar h$; one more use of Tonelli's theorem gives $\int_{\R^{n-1}}\bar h=\int_{\R^n}h$.
\end{proof}

Call $p\colon\R^n\to[0,\infty)$ \term{box-regular} if, for some bounded open box $Q$, it is continuous on the closed box $\overline Q$, positive on $Q$, and zero outside $\overline Q$.

\begin{lemma}[Knothe's map]\label{lem:transport:knothe}
Let $p$ and $q$ be box-regular probability densities, positive on the boxes $Q=\prod_k(a_k,b_k)$ and $Q'=\prod_k(a'_k,b'_k)$. There is a Borel map $T\colon Q\to Q'$ such that each $T_k$ depends only on $x_{\le k}$, each function $s\mapsto T_k(x_{<k},s)$ is a continuously differentiable increasing bijection of $(a_k,b_k)$ onto $(a'_k,b'_k)$, and
\begin{equation}\label{eq:transport:knothejac}
\prod_{k=1}^n\partial_kT_k(x)=\frac{p(x)}{q(T(x))}\qquad\text{for all }x\in Q.
\end{equation}
\end{lemma}

\begin{proof}
The plan is to rearrange the conditional law of each coordinate given the previous ones. For $0\le k\le n$ let $p_k(x_{\le k})$ be the integral of $p$ over the last $n-k$ coordinates, so that $p_0=1$ and $p_n=p$. Uniform continuity of $p$ on $\overline Q$ makes each $p_k$ continuous on $\prod_{j\le k}[a_j,b_j]$, and positivity of $p$ on $Q$ makes it positive on $\prod_{j\le k}(a_j,b_j)$. For $x\in Q$ put
\[
F_k(x_{\le k})=\frac{1}{p_{k-1}(x_{<k})}\int_{a_k}^{x_k}p_k(x_{<k},s)\dd s.
\]
By Fubini's theorem $\int_{a_k}^{b_k}p_k(x_{<k},s)\dd s=p_{k-1}(x_{<k})$, so for fixed $x_{<k}$ the function $F_k$ is a continuously differentiable increasing bijection of $(a_k,b_k)$ onto $(0,1)$ with derivative $p_k/p_{k-1}$; moreover $F_k$ is jointly continuous. Define $G_k$ from $q$ in the same way, and define $T_1,\dots,T_n$ recursively by $G_k\bigl(T_1(x),\dots,T_k(x)\bigr)=F_k(x_{\le k})$. By the inverse function theorem in one variable, $s\mapsto T_k(x_{<k},s)$ is a continuously differentiable increasing bijection of $(a_k,b_k)$ onto $(a'_k,b'_k)$, and
\[
\partial_kT_k(x)=\frac{p_k(x_{\le k})\,/\,p_{k-1}(x_{<k})}{q_k(T_{\le k}(x))\,/\,q_{k-1}(T_{<k}(x))}.
\]
The product over $k$ telescopes to $p_n(x)/q_n(T(x))$, which is \eqref{eq:transport:knothejac}. Finally $T_k$ is Borel, by induction on $k$, because $T_k(x)$ is the supremum of the rational $s\in(a'_k,b'_k)$ with $G_k(T_{<k}(x),s)<F_k(x_{\le k})$.
\end{proof}

One can check that $T$ transports $p\dd x$ to $q\dd x$; we need only \eqref{eq:transport:knothejac}.

\begin{theorem}[Pr\'ekopa--Leindler for box-regular functions]\label{thm:transport:plbox}
Let $0<t<1$, let $f$ and $g$ be box-regular, and let $h\colon\R^n\to[0,\infty)$ be a Borel function satisfying \eqref{eq:transport:plhyp}. Then $\int h\ge\bigl(\int f\bigr)^{1-t}\bigl(\int g\bigr)^t$.
\end{theorem}

\begin{proof}
The integrals $F=\int f$ and $G=\int g$ are positive. Let $T$ be the map of \cref{lem:transport:knothe} for the densities $f/F$ and $g/G$, and let $z(x)=(1-t)x+tT(x)$ on the box $Q$ of $f$. The map $z$ is triangular, and $\partial_kz_k=(1-t)+t\,\partial_kT_k$ is positive and continuous in $x_k$, so \cref{lem:transport:triangular} gives
\[
\int h\ \ge\ \int_Q h(z(x))\prod_{k=1}^n\bigl((1-t)+t\,\partial_kT_k(x)\bigr)\dd x.
\]
We bound the integrand pointwise. The hypothesis \eqref{eq:transport:plhyp} with $y=T(x)$ gives $h(z(x))\ge f(x)^{1-t}g(T(x))^t$, where $g(T(x))>0$ because $T(x)$ lies in the box of $g$. The inequality $(1-t)+t\lambda\ge\lambda^t$ for $\lambda>0$, applied to each factor, and \eqref{eq:transport:knothejac} give
\[
\prod_{k=1}^n\bigl((1-t)+t\,\partial_kT_k(x)\bigr)\ \ge\ \Bigl(\prod_{k=1}^n\partial_kT_k(x)\Bigr)^{t}=\Bigl(\frac{G\,f(x)}{F\,g(T(x))}\Bigr)^{t}.
\]
So the integrand is at least $(G/F)^tf(x)$, and integrating over $Q$ gives $\int h\ge(G/F)^tF=F^{1-t}G^t$.
\end{proof}

Compared with the classical proof, the induction on dimension has disappeared into the triangular Jacobian, and the level sets into a map that matches mass instead of height.

\begin{corollary}\label{cor:transport:bm}
The Brunn--Minkowski inequality \eqref{eq:transport:bm} holds for all nonempty compact $A,B\subseteq\R^n$.
\end{corollary}

\begin{proof}
If $\abs{A}=0$, then $A+B$ contains a translate of $B$ and \eqref{eq:transport:bm} holds; likewise if $\abs{B}=0$. Fix $t,\eps,\delta\in(0,1)$, put $m=\min\{t,1-t\}$, let $A^\eps$ and $B^\eps$ be the closed $\eps$-neighborhoods of $A$ and $B$, and let $Q$ be a bounded open box containing $A^1\cup B^1$. The functions $f=\max\{\delta,1-\dist(\cdot,A)/\eps\}$ and $g=\max\{\delta,1-\dist(\cdot,B)/\eps\}$ on $\overline Q$, extended by $0$, are box-regular, with $f\ge\ind_A$ and $g\ge\ind_B$. Let $h=1$ on $K=(1-t)A^\eps+tB^\eps$, $h=\delta^m$ on $\overline Q\setminus K$, and $h=0$ elsewhere. Since $f^{1-t}g^t\le1$, with $f(x)^{1-t}g(y)^t\le\delta^m$ unless $x\in A^\eps$ and $y\in B^\eps$, and since $\overline Q$ is convex, $h$ satisfies \eqref{eq:transport:plhyp}, and \cref{thm:transport:plbox} gives $\abs{K}+\delta^m\abs{Q}\ge\abs{A}^{1-t}\abs{B}^t$. Let $\delta\to0$ and then $\eps\to0$; by compactness the sets $K$ decrease to $(1-t)A+tB$, so $\abs{(1-t)A+tB}\ge\abs{A}^{1-t}\abs{B}^t$. With $a=\abs{A}^{1/n}$, $b=\abs{B}^{1/n}$, and $t=b/(a+b)$, apply this to the unit-volume sets $A/a$ and $B/b$, whose combination $(1-t)A/a+tB/b$ is $(A+B)/(a+b)$; the result is $\abs{A+B}\ge(a+b)^n$.
\end{proof}

For compact $A$ and the closed unit ball $B$, the outer \term{Minkowski content}\index{Minkowski content} of $A$ is $\liminf_{\eps\to0^+}(\abs{A+\eps B}-\abs{A})/\eps$; for a set with smooth boundary it equals the area of the boundary. Let $\omega_n=\abs{B}$.

\begin{corollary}[Isoperimetric inequality]\label{cor:transport:iso}
Every compact set $A\subseteq\R^n$ with $\abs{A}>0$ satisfies
\[
\liminf_{\eps\to0^+}\frac{\abs{A+\eps B}-\abs{A}}{\eps}\ \ge\ n\,\omega_n^{1/n}\abs{A}^{(n-1)/n},
\]
with equality when $A$ is a ball.
\end{corollary}

\begin{proof}
Apply \cref{cor:transport:bm} to $A$ and $\eps B$, and keep two terms of the binomial expansion:
\[
\abs{A+\eps B}\ \ge\ \bigl(\abs{A}^{1/n}+\eps\,\omega_n^{1/n}\bigr)^n\ \ge\ \abs{A}+n\eps\,\omega_n^{1/n}\abs{A}^{(n-1)/n}.
\]
Subtract $\abs{A}$, divide by $\eps$, and let $\eps\to0^+$. For $A=rB$, $\abs{A+\eps B}=\omega_n(r+\eps)^n$, and the limit is $n\omega_nr^{n-1}=n\,\omega_n^{1/n}\abs{A}^{(n-1)/n}$.
\end{proof}

No minimizer is needed, so the objection of Weierstrass does not arise. Minkowski knew how to pass from Brunn--Minkowski to isoperimetry for convex bodies; transport supplies a proof of Brunn--Minkowski that is the same as every other proof here, and Gromov's direct proof of isoperimetry (Exercise~\ref{ex:transport:gromov}).

\subsection{The Gaussian inequalities, and every uniformly log-concave measure}\label{sec:transport:gaussafter}

Let $V\colon\R\to\R$ be continuously differentiable and $\rho$-convex, $\rho>0$. Then $e^{-V}$ is bounded by a multiple of a Gaussian density, so $\pi=Z^{-1}e^{-V(x)}\dd x$ is a probability measure with moments of all orders; the standard Gaussian is $V(x)=x^2/2$, $\rho=1$. Let $\transportReg_\pi$ be the set of probability measures $\mu=f\pi$ whose density $f$ is continuously differentiable, positive, and equal to $1$ outside a bounded interval. For such $\mu$, $\transportH{\mu}{\pi}=\int f\log f\dd\pi$ and $\transportI{\mu}{\pi}=\int f'^2/f\dd\pi$ are finite.

\begin{lemma}[The monotone map between regular measures]\label{lem:transport:monotone}
Let $\mu=f\pi$ and $\nu=g\pi$ belong to $\transportReg_\pi$, with distribution functions $F$ and $G$. The map $T=G^{-1}\circ F$ is a continuously differentiable increasing bijection of $\R$ that transports $\mu$ to $\nu$ optimally for the quadratic cost, equals the identity outside a bounded interval, and satisfies
\begin{equation}\label{eq:transport:jacobian1}
f(x)\,e^{-V(x)}=g(T(x))\,e^{-V(T(x))}\,T'(x)\qquad\text{for all }x\in\R.
\end{equation}
\end{lemma}

\begin{proof}
The function $F$ has the positive continuous derivative $Z^{-1}fe^{-V}$ and increases from $0$ to $1$, so it is a continuously differentiable increasing bijection of $\R$ onto $(0,1)$; so is $G$. Hence $T$ is a continuously differentiable increasing bijection of $\R$, and differentiating $G(T(x))=F(x)$ gives \eqref{eq:transport:jacobian1}. By \cref{thm:transport:oned}, $T$ transports $\mu$ to $\nu$ optimally; both measures have finite second moments because $\pi$ does. Suppose that $f=g=1$ outside $[-R,R]$, and let $\Pi$ be the distribution function of $\pi$. For $x<-R$ we have $F(x)=\Pi(x)<\Pi(-R)=G(-R)$, so $T(x)<-R$, where $G=\Pi$; then $G(T(x))=F(x)$ reads $\Pi(T(x))=\Pi(x)$, so $T(x)=x$. The same argument with $1-F$ and $1-G$ gives $T(x)=x$ for $x>R$.
\end{proof}

\begin{theorem}[One inequality, three consequences]\label{thm:transport:master}
Let $\mu=f\pi$ and $\nu=g\pi$ belong to $\transportReg_\pi$, let $T$ be the monotone map of \cref{lem:transport:monotone}, and let $u(x)=T(x)-x$. Then
\begin{equation}\label{eq:transport:master}
\transportH{\nu}{\pi}\ \ge\ \transportH{\mu}{\pi}+\int u\,(\log f)'\dd\mu+\frac{\rho}{2}\int u^2\dd\mu.
\end{equation}
Consequently every $\mu$ and $\nu$ in $\transportReg_\pi$ satisfy:
\begin{enumerate}
\item \emph{(Talagrand)} $\transportW(\nu,\pi)^2\le\frac{2}{\rho}\,\transportH{\nu}{\pi}$;
\item \emph{(HWI)} $\transportH{\mu}{\pi}\le\transportW(\mu,\pi)\sqrt{\transportI{\mu}{\pi}}-\frac{\rho}{2}\,\transportW(\mu,\pi)^2$;
\item \emph{(logarithmic Sobolev)} $\transportH{\mu}{\pi}\le\frac{1}{2\rho}\,\transportI{\mu}{\pi}$.
\end{enumerate}
\end{theorem}

\begin{proof}
The plan is to take logarithms in \eqref{eq:transport:jacobian1}, bound the two nonlinear terms by convexity, integrate against $\mu$, and integrate by parts once.

\emph{The pointwise inequality.} Taking logarithms in \eqref{eq:transport:jacobian1} gives
\[
\log g(T(x))=\log f(x)+\bigl[V(T(x))-V(x)\bigr]-\log T'(x).
\]
The $\rho$-convexity of $V$ bounds the bracket below by $V'(x)u(x)+\frac{\rho}{2}u(x)^2$, and $\log T'(x)=\log(1+u'(x))\le u'(x)$ because $T'>0$. Hence
\[
\log g(T(x))\ \ge\ \log f(x)+V'(x)u(x)-u'(x)+\frac{\rho}{2}u(x)^2.
\]

\emph{Integration.} Both sides are continuous and vanish outside a bounded interval, where $u=0$ and $f=g\circ T=1$. Integrate against $\mu$. Because $T$ transports $\mu$ to $\nu$, the left side becomes $\int\log g\dd\nu=\transportH{\nu}{\pi}$, and the first term on the right becomes $\transportH{\mu}{\pi}$. The function $u\,f\,e^{-V}$ is continuously differentiable and vanishes outside a bounded interval, so its derivative integrates to zero: $\int(u'-V'u)f\dd\pi+\int uf'\dd\pi=0$. Therefore $\int(V'u-u')\dd\mu=\int uf'\dd\pi=\int u\,(\log f)'\dd\mu$, which proves \eqref{eq:transport:master}.

\emph{Talagrand.} Take $\mu=\pi$, that is, $f=1$. Then \eqref{eq:transport:master} and \eqref{eq:transport:mapbound} give $\transportH{\nu}{\pi}\ge\frac{\rho}{2}\int(T(x)-x)^2\dd\pi(x)\ge\frac{\rho}{2}\transportW(\pi,\nu)^2$.

\emph{HWI.} Take $\nu=\pi$, that is, $g=1$, and apply the Cauchy--Schwarz inequality to the middle term of \eqref{eq:transport:master}. With $\norm{u}^2=\int u^2\dd\mu$,
\[
\transportH{\mu}{\pi}\ \le\ -\int u\,(\log f)'\dd\mu-\frac{\rho}{2}\norm{u}^2\ \le\ \norm{u}\sqrt{\transportI{\mu}{\pi}}-\frac{\rho}{2}\norm{u}^2,
\]
and $\norm{u}=\transportW(\mu,\pi)$ because the monotone map is optimal (\cref{lem:transport:monotone}).

\emph{Logarithmic Sobolev.} Complete the square: for $w,I\ge0$, $w\sqrt I-\frac{\rho}{2}w^2=\frac{I}{2\rho}-\frac{1}{2\rho}\bigl(\sqrt I-\rho w\bigr)^2\le\frac{I}{2\rho}$. Apply this to HWI.
\end{proof}

\index{HWI inequality}Set this beside Gross's route: no cube, no central limit theorem, no semigroup, only the logarithm of a derivative and one integration by parts, and no optimality of $T$ except in HWI. Nor is product structure used: the Gaussian is one $\rho$-convex potential among many, and the same lines prove \eqref{eq:transport:be} on the line. Three inequalities come out of one, with one constant, and translations ($T'=1$) make each step an equality for the Gaussian, which is why its two inequalities are sharp for the same measures. The restriction to $\transportReg_\pi$ is removed by approximation \parencite{Villani2003}, and in $\R^n$ the computation runs with Brenier's map.

\begin{theorem}[Transport inequalities in $\R^n$]\label{thm:transport:ndim}
Let $V\colon\R^n\to\R$ be twice continuously differentiable with $D^2V\ge\rho I$, $\rho>0$, and let $\pi=Z^{-1}e^{-V(x)}\dd x$. Then Talagrand's inequality $\transportW(\nu,\pi)^2\le\frac{2}{\rho}\transportH{\nu}{\pi}$ holds for every probability measure $\nu$ on $\R^n$, and HWI and the logarithmic Sobolev inequality, in the form of \cref{thm:transport:master}, hold for every $\nu=f\pi$ with $f$ smooth and compactly supported, with $\transportI{\nu}{\pi}=\int_{\{f>0\}}\abs{\nabla f}^2/f\dd\pi$.
\end{theorem}

\begin{proofsketch}
Repeat the proof of \cref{thm:transport:master} with the map $x+\nabla\theta(x)$ of \cref{thm:transport:brenier}: the Monge--Amp\`ere equation replaces \eqref{eq:transport:jacobian1}, and $\log\det(I+D^2\theta)\le\Delta\theta$ replaces $\log(1+u')\le u'$ \parencite{CorderoErausquin2002}. Omitted is regularity: $D^2\theta$ exists only almost everywhere, and the integration by parts becomes an inequality, in the right direction because the distributional Laplacian of $\theta$ dominates its almost-everywhere Laplacian, $\frac12\abs{x}^2+\theta$ being convex.
\end{proofsketch}

For $\pi=\gamma$ this gives \cref{thm:transport:gross,thm:transport:talagrand}; for general $V$ it gives \eqref{eq:transport:be} and the HWI inequality of \textcite{OttoVillani2000}.

\section{After: the geometry behind the inequalities}\label{sec:transport:after2}

One inequality, \eqref{eq:transport:master}, produced three. The geometry explains why.

\subsection{Entropy is convex along straight lines}\label{sec:transport:dispconvex}

\begin{theorem}[Displacement convexity]\label{thm:transport:dispconvex}
If $V\colon\R^n\to\R$ is $\rho$-convex, $\rho>0$, and $\pi=Z^{-1}e^{-V(x)}\dd x$, then $\transportH{\cdot}{\pi}$ is displacement $\rho$-convex on the absolutely continuous measures in $\cP_2(\R^n)$.
\end{theorem}

\begin{proofsketch}
For $\mu=\varrho\dd x$, $\transportH{\mu}{\pi}=\int\varrho\log\varrho\dd x+\int V\dd\mu+\log Z$. The potential term inherits $\rho$-convexity from $V$ along the segments from $x$ to $T(x)$. For the entropy, \textcite{McCann1997} wrote $\int\varrho_t\log\varrho_t=\int\varrho_0\log\varrho_0-\int\log\det((1-t)I+tD^2\varphi)\dd\mu_0$, convex in $t$ because $\log\det$ is concave; Exercise~\ref{ex:transport:mccann1d} does this rigorously on the line.
\end{proofsketch}

\index{displacement convexity}Now read \eqref{eq:transport:master} again. Along the displacement interpolation from $\mu$ to $\nu$, the $\rho$-convex function $t\mapsto\transportH{\mu_t}{\pi}$ lies above its tangent at $t=0$ plus $\frac{\rho}{2}t^2\transportW(\mu,\nu)^2$, and formally its derivative at $t=0$ is $\int u\,(\log f)'\dd\mu$. So \eqref{eq:transport:master} is the first-order form of displacement convexity, and the three inequalities are one convexity statement read at $\mu=\pi$ and at $\nu=\pi$.

\subsection{A formal Riemannian geometry}\label{sec:transport:otto}

This subsection is \emph{formal}: it treats $\cP_2(\R^n)$ as a smooth Riemannian manifold, following the heuristics of \textcite{Otto2001}, now called \term{Otto calculus};\index{Otto calculus} the rigorous theory is in \textcite{AmbrosioGigliSavare2008}. A curve $\varrho_t$ moves with velocity $v_t$ if $\partial_t\varrho_t+\nabla\cdot(\varrho_tv_t)=0$, and $v_t$ has squared length $\int\abs{v_t}^2\dd\varrho_t$; by the Benamou--Brenier theorem the resulting length metric is $\transportW$. An integration by parts gives
\[
\frac{\dd}{\dd t}\transportH{\varrho_t}{\pi}=\int\nabla\log\frac{\varrho_t}{\pi}\cdot v_t\dd\varrho_t,
\]
so the gradient of $\transportH{\cdot}{\pi}$ at $\varrho$ is $\nabla\log(\varrho/\pi)$, whose squared length is $\transportI{\varrho}{\pi}$, and the gradient flow is the \term{Fokker--Planck equation}\index{Fokker--Planck equation}
\[
\partial_t\varrho=\nabla\cdot\bigl(\varrho\nabla\log(\varrho/\pi)\bigr)=\Delta\varrho+\nabla\cdot(\varrho\nabla V)
\]
of the Langevin diffusion $\dd X_t=-\nabla V(X_t)\dd t+\sqrt2\dd B_t$. \textcite{JordanKinderlehrerOtto1998} made this rigorous: the scheme $\varrho_{k+1}=\argmin_\varrho\{\transportH{\varrho}{\pi}+\transportW(\varrho,\varrho_k)^2/(2\tau)\}$ converges to the Fokker--Planck flow as $\tau\to0$. Displacement $\rho$-convexity bounds the Hessian below, and \cref{tab:transport:dictionary} completes the translation.

\begin{table}[t]
\centering
\caption{A $\rho$-convex function $f$ on $\R^d$ and the relative entropy on Wasserstein space. Read $f(x)-f^\star$ as $\transportH{\nu}{\pi}$, $\abs{\nabla f(x)}$ as $\sqrt{\transportI{\nu}{\pi}}$, and $\abs{x-x^\star}$ as $\transportW(\nu,\pi)$.}
\label{tab:transport:dictionary}
\small
\begin{tabular}{@{}ll@{}}
\toprule
Euclidean space & Wasserstein space\\
\midrule
$f-f^\star\le\abs{\nabla f}\abs{x-x^\star}-\frac\rho2\abs{x-x^\star}^2$ & HWI\\
$f-f^\star\le\abs{\nabla f}^2/(2\rho)$ & logarithmic Sobolev\\
$f-f^\star\ge\frac\rho2\abs{x-x^\star}^2$ & Talagrand\\
$f(x_t)-f^\star\le e^{-2\rho t}(f(x_0)-f^\star)$ & decay along Fokker--Planck\\
domination implies growth & Otto--Villani theorem\\
\bottomrule
\end{tabular}
\end{table}

In optimization, gradient domination is the \term{Polyak--{\L}ojasiewicz inequality}\index{Polyak--{\L}ojasiewicz inequality} (Polyak 1963). For a $\rho$-convex $f$ the first four rows are elementary. The last needs no convexity: along the gradient flow from $x$, $\sqrt{f(x_t)-f^\star}$ decreases at least $\sqrt{\rho/2}$ times as fast as the path grows, so the path reaches a minimizer within distance $\sqrt{2(f(x)-f^\star)/\rho}$ (Exercise~\ref{ex:transport:pl}).

\begin{theorem}[\textcite{OttoVillani2000}]\label{thm:transport:ottovillani}
Let $\pi$ be a probability measure on $\R^n$ such that $\transportH{\nu}{\pi}\le\frac{1}{2\rho}\transportI{\nu}{\pi}$ for every $\nu=f\pi$ with $f$ smooth and positive. Then $\transportW(\nu,\pi)^2\le\frac{2}{\rho}\transportH{\nu}{\pi}$ for every probability measure $\nu$ on $\R^n$.
\end{theorem}

\begin{proofsketch}[Formal proof]
Repeat the Euclidean argument along the Fokker--Planck flow from $\nu$: the entropy decreases at rate $\transportI{\varrho_t}{\pi}$, the metric speed is $\sqrt{\transportI{\varrho_t}{\pi}}$, and the logarithmic Sobolev inequality bounds the length of the curve from $\nu$ to $\pi$, hence $\transportW(\nu,\pi)$, by $\sqrt{2\transportH{\nu}{\pi}/\rho}$. Otto and Villani made this rigorous; Bobkov, Gentil, and Ledoux later gave a short proof through the Hamilton--Jacobi equation.
\end{proofsketch}

Gradient domination implies quadratic growth, so the logarithmic Sobolev inequality implies Talagrand's, even for measures that are not log-concave.

\subsection{Dynamics: sampling and wide networks}\label{sec:transport:dynamics}

For the law $\varrho_t$ of the \term{Langevin dynamics},\index{Langevin dynamics} the dissipation identity $\frac{\dd}{\dd t}\transportH{\varrho_t}{\pi}=-\transportI{\varrho_t}{\pi}$, formal here and rigorous for regular initial laws, and a logarithmic Sobolev inequality give $\transportH{\varrho_t}{\pi}\le e^{-2\rho t}\transportH{\varrho_0}{\pi}$; for $V(x)=x^2/2$ and $\varrho_0=N(m,1)$, $\varrho_t=N(me^{-t},1)$ attains the rate. The discretization $X_{k+1}=X_k-\eta\nabla V(X_k)+\sqrt{2\eta}\,\xi_k$, with independent standard Gaussian $\xi_k$, samples from $\pi$ approximately; Vempala and Wibisono (2019) showed that a logarithmic Sobolev inequality and a Lipschitz gradient, without convexity, make it converge in relative entropy up to a bias proportional to $\eta$.

A two-layer network $x\mapsto\frac1m\sum_i\phi(x;w_i)$ depends on its parameters only through their empirical measure, which gradient descent moves by a continuity equation (Exercise~\ref{ex:transport:meanfield}); as $m\to\infty$, training becomes a Wasserstein gradient flow of the risk \parencite{ChizatBach2018,MeiMontanariNguyen2018}. The risk is convex along mixtures, though not in the parameters, and under homogeneity and support conditions Chizat and Bach showed that a limit of the flow is a global minimizer.

\paragraph{New reach.}\label{sec:transport:reach} The proofs use only a map and a pointwise inequality, so they travel. Riemannian manifolds have no boxes but have optimal maps, and Cordero-Erausquin, McCann, and Schmuckenschl\"ager (2001) proved Pr\'ekopa--Leindler inequalities there, with Ricci curvature in the Jacobian; Sturm (2006) and \textcite{LottVillani2009} then \emph{defined} lower Ricci bounds for metric measure spaces by displacement convexity of entropy. In $\R^n$, transport gave sharp Sobolev inequalities with their extremals (Cordero-Erausquin, Nazaret, and Villani 2004), the reverse Brascamp--Lieb inequality (Barthe 1998), and sharp quantitative anisotropic isoperimetry (Figalli, Maggi, and Pratelli 2010).

\section{Costs and limits}\label{sec:transport:costs}

\paragraph{Where did the difficulty go?}
Into the maps. Here it was paid in three lemmas (the monotone map, Knothe's map, the triangular change of variables) and one approximation. With Brenier's map it is paid in Brenier's theorem (Kantorovich duality and convex analysis), in the almost-everywhere Monge--Amp\`ere equation (Aleksandrov's theorem), and, for smooth maps, in Caffarelli's regularity theory, which needs among other things a convex target. The geometric half is paid in the metric theory of gradient flows \parencite{AmbrosioGigliSavare2008}. The price is paid once: the same maps prove every inequality of the chapter, each at the cost of a pointwise inequality and an integration by parts. For Brunn--Minkowski alone, the transport proof is not shorter than Hadwiger and Ohmann's once the approximation is counted; there the dividend is unification and reach.

\paragraph{What did not dissolve.}
Convexity is the fuel. For a double-well potential \eqref{eq:transport:master} says nothing, and the Kannan--Lov\'asz--Simonovits problem on log-concave measures without uniform convexity needed other tools, notably stochastic localization. Transport inequalities give concentration (Exercise~\ref{ex:transport:marton}) but not sharp isoperimetry, which in Gauss space is the stronger theorem of Borell and of Sudakov and Tsirelson. In statistics, transport suffers a curse of dimensionality: the empirical measure of $n$ samples from an absolutely continuous law on $[0,1]^d$, $d\ge3$, is at $W_1$-distance of order $n^{-1/d}$ from the law (Dudley 1969; Fournier and Guillin 2015). Optimal couplings of $n$ points solve an assignment problem; entropic regularization \parencite{Cuturi2013} is fast but changes the object. Monge's own cost $\abs{x-y}$ was the hardest case, settled only around 2000 \parencite{Villani2009}.

\paragraph{Other seas.}
The semigroup approach of Bakry and \'Emery is a theory of its own, often sharper about dimension; stochastic proofs build couplings from Brownian motion instead of maps.

\section{The move, distilled}\label{sec:transport:distilled}

\begin{move}{Geometrization}
\movepart{In one sentence} Carry one object onto the other by a well-chosen map, so that the inequality becomes a pointwise inequality for the Jacobian, integrated once; then read the result as convexity along the straight lines that the maps draw in the space of measures.
\movepart{Triggers} The two sides are integrals or volumes of two objects of the same kind; translates or dilates are extremal; the known proofs induct on dimension or tensorize; the constant is dimension-free; a functional of measures should be convex but is not convex along mixtures; a dynamics dissipates an entropy.
\movepart{Recipe} Choose the map: monotone on the line, Knothe's for elementary proofs, Brenier's for optimality or a symmetric Jacobian. Write the Jacobian equation, apply a pointwise inequality (arithmetic--geometric mean, concavity of $\log\det$, $\log x\le x-1$, convexity of the potential), and integrate by parts. Read the result as displacement convexity, let Euclidean convex analysis suggest the next statement, then pay for rigor.
\movepart{Anti-pattern} Calling an integration by parts a Wasserstein gradient flow when the geometry adds nothing; invoking a map whose existence or regularity is as hard as the inequality; estimating transport distances in high dimension as if at parametric rates.
\end{move}

\dissolutiontest{Gaussian and Bakry--\'Emery logarithmic Sobolev inequalities on the line in half a page; isoperimetry without a minimizer; Pr\'ekopa--Leindler in all dimensions at once.}{No product structure, so every uniformly log-concave measure; no existence of minimizers; no boxes, so manifolds are allowed.}{Brunn--Minkowski, Pr\'ekopa--Leindler, isoperimetry, and the logarithmic Sobolev, Talagrand, and HWI inequalities are one computation.}{HWI and Otto--Villani; diffusions as gradient flows; synthetic Ricci curvature; sharp Sobolev inequalities; mean-field limits of wide networks.}

Kantorovich's problem is a linear program whose dual variables are the multipliers of \cref{ch:lpduality}, and Brenier's potential is assembled from them. The geometric half is a cousin of \cref{ch:regret}, where dynamics replaced an existence theorem: the Otto--Villani theorem follows a gradient flow and measures its length. Existence of optimal couplings rests on compactness of sets of measures, as in the completion move of \cref{ch:lebesgue}. On the discrete cube, where Gross began, Fourier analysis does the work instead (\cref{ch:fourier}): transport needs a continuum and convexity, Fourier analysis a group.

\begin{exercises}
\item \consolidation\label{ex:transport:gauss} \emph{Gaussian equality cases.} (a) Use \eqref{eq:transport:w2line} to show that $\transportW(N(m_1,s_1^2),N(m_2,s_2^2))^2=(m_1-m_2)^2+(s_1-s_2)^2$. (b) Show that $\gamma_m=N(m,1)$ has $\transportH{\gamma_m}{\gamma_0}=m^2/2$ and $\transportI{\gamma_m}{\gamma_0}=m^2$, and is an equality case of all three Gaussian inequalities.

\item \consolidation\label{ex:transport:convexcost} \emph{Convex costs.} Let $x_1\le\dots\le x_m$, $y_1\le\dots\le y_m$, and $h\colon\R\to\R$ convex. Show that the identity minimizes $\sum_ih(x_i-y_{\sigma(i)})$ over permutations $\sigma$, so the coupling of \cref{thm:transport:oned} is optimal for the cost $h(x-y)$ between uniform measures on $m$ points (use the Birkhoff--von Neumann theorem of \cref{ch:lpduality}). \hint{for $x\le x'$ and $y\le y'$, convexity gives $h(x-y)+h(x'-y')\le h(x-y')+h(x'-y)$; uncross pairs one at a time.}

\item \consolidation\label{ex:transport:knothe} \emph{Knothe is not Brenier.} Let $A=(0,1)^2$ and $B=\{(y_1,y_2):0<y_2<y_1<1\}$ carry uniform probability measures. Show that Knothe's map from $A$ to $B$ is $T(x)=(\sqrt{x_1},\,x_2\sqrt{x_1})$, check that $\partial_1T_1\,\partial_2T_2=\abs{B}/\abs{A}$, and show that $T$ is not the gradient of a convex function, hence not optimal (\cref{thm:transport:brenier}). \hint{a continuously differentiable gradient has a symmetric Jacobian matrix.}

\item \consolidation\label{ex:transport:mccann1d} \emph{Displacement convexity on the line.} Let $f$ and $g$ be probability densities, continuous on bounded closed intervals, positive inside, and zero outside; let $T$ be the monotone map from $f\dd x$ to $g\dd x$, and $\mu_t=((1-t)\mathrm{id}+tT)\transportpf(f\dd x)$. (a) Show that $\mu_t$ has a density $\varrho_t$ with $\varrho_t((1-t)x+tT(x))=f(x)/((1-t)+tT'(x))$, and that $t\mapsto\int\varrho_t\log\varrho_t$ is convex. (b) Show that the variance of $\mu_t$ is convex in $t$, and that of $(1-t)\mu_0+t\mu_1$ concave.

\item \consolidation\hard\label{ex:transport:pl} \emph{Gradient domination implies quadratic growth.} Let $f\colon\R^d\to\R$ have a locally Lipschitz gradient, attain its minimum $f^\star$ on a set $M$, and satisfy $f(x)-f^\star\le\abs{\nabla f(x)}^2/(2\rho)$ for all $x$, with $\rho>0$. Prove that $f(x)-f^\star\ge\frac\rho2\dist(x,M)^2$. \hint{along $\dot x_t=-\nabla f(x_t)$, $\frac{\dd}{\dd t}\sqrt{f(x_t)-f^\star}\le-\sqrt{\rho/2}\,\abs{\dot x_t}$ while $f(x_t)>f^\star$; treat separately the case in which the flow reaches $M$ in finite time.}

\item \transfer\hard\label{ex:transport:marton} \emph{Marton's argument.} Assume that $\transportW(\nu,\gamma_n)^2\le2\transportH{\nu}{\gamma_n}$ for every $\nu$, with $\gamma_n$ the standard Gaussian measure (\cref{thm:transport:ndim}). For Borel $A$ with $\gamma_n(A)\ge\frac12$ and $A_r=\{x:\dist(x,A)\le r\}$, prove that $\gamma_n(A_r)\ge1-\exp(-(r-\sqrt{2\log2})^2/2)$ for $r\ge\sqrt{2\log2}$, in every dimension. Deduce that a $1$-Lipschitz $F$ with median $m_F$ satisfies $\gamma_n(F>m_F+r+\sqrt{2\log2})\le e^{-r^2/2}$ for $r\ge0$. \hint{compare $\gamma_n$ conditioned on $A$ and on the complement of $A_r$ through the triangle inequality for $\transportW$.}

\item \transfer\hard\label{ex:transport:gromov} \emph{Gromov's isoperimetric proof.} Let $\Omega\subseteq\R^n$ be a bounded open set with $C^1$ boundary, $B$ the open unit ball, and $\omega_n=\abs{B}$. Let $T$ be continuously differentiable near $\overline\Omega$, with $T(\overline\Omega)\subseteq\overline B$, triangular with $\partial_kT_k>0$ and $\prod_k\partial_kT_k=\omega_n/\abs{\Omega}$ on $\Omega$ (as is Knothe's map from $\Omega$ to $B$ when regular enough). Prove that $\partial\Omega$ has area at least $n\,\omega_n^{1/n}\abs{\Omega}^{(n-1)/n}$. \hint{$(\omega_n/\abs{\Omega})^{1/n}\le\frac1n\operatorname{div}T$; integrate, and use the divergence theorem and $\abs{T}\le1$.}

\item \transfer\label{ex:transport:meanfield} \emph{Wide networks.} Let $\phi\colon\R^d\times\R^p\to\R$ be smooth and bounded with bounded derivatives, and $(X,Y)$ random with $\E Y^2<\infty$. For a probability measure $\varrho$ on $\R^p$ let $f_\varrho=\int\phi(\cdot;w)\dd\varrho(w)$, $R(\varrho)=\frac12\E(f_\varrho(X)-Y)^2$, and $\Psi_\varrho(w)=\E[(f_\varrho(X)-Y)\phi(X;w)]$. Let $\dot w_i=-m\nabla_{w_i}R(\varrho_t)$, with $\varrho_t=\frac1m\sum_{i=1}^m\delta_{w_i(t)}$. Show that $\frac{\dd}{\dd t}\int\zeta\dd\varrho_t=-\int\nabla\zeta\cdot\nabla\Psi_{\varrho_t}\dd\varrho_t$ for smooth compactly supported $\zeta$, and $\frac{\dd}{\dd t}R(\varrho_t)=-\int\abs{\nabla\Psi_{\varrho_t}}^2\dd\varrho_t$.

\item \research\harder\label{ex:transport:generalization} \emph{Transport and generalization.} Let $S=(Z_1,\dots,Z_n)$ be independent with law $\mu$ on $\R^k$, let $W$ be the output of a randomized learning algorithm on $S$, and let the loss $\ell(w,z)$ be $L$-Lipschitz in $z$. (a) Suppose that $W_1(\nu,\mu)\le\sqrt{2C\transportH{\nu}{\mu}}$ for every $\nu$; show that $\E\exp(\lambda(\ell(w,Z_1)-\E\ell(w,Z_1)))\le\exp(\lambda^2L^2C/2)$ for all $w$ and real $\lambda$. (b) Deduce from the mutual-information bound of \textcite{XuRaginsky2017} that the expected generalization gap is at most $\sqrt{2L^2C\,I(W;S)/n}$. (c) Deterministic algorithms such as gradient descent usually have $I(W;S)=\infty$; find a transport-based substitute for $I(W;S)$ that is finite for them and still controls the gap. This part is open-ended. \hint{for (a), combine $\int F\dd\nu-\int F\dd\mu\le W_1(\nu,\mu)$ for $1$-Lipschitz $F$ with $\log\int e^{G}\dd\mu=\sup_\nu(\int G\dd\nu-\transportH{\nu}{\mu})$ for bounded $G$.}
\end{exercises}

\notesandrefs

The transport problem goes back to \textcite{Monge1781}, its relaxation and duality to \textcite{Kantorovich1942}. Brenier's theorem is from \textcite{Brenier1991}; \textcite{McCann1995} removed the moment assumptions, and \textcite{McCann1997} introduced displacement interpolation and convexity. The standard references are \textcite{Villani2003,Villani2009}, and \textcite{Gardner2002} is the guide to Brunn--Minkowski. Gromov used Knothe's map for isoperimetry in an appendix to the 1986 lecture notes of Milman and Schechtman.

The logarithmic Sobolev inequality is due to \textcite{Gross1975}, the curvature criterion to \textcite{BakryEmery1985}, and the Gaussian transport inequality to \textcite{Talagrand1996}; \textcite{OttoVillani2000} linked them through HWI, and \textcite{CorderoErausquin2002} gave the transport proofs followed here. The geometric side is due to \textcite{JordanKinderlehrerOtto1998} and \textcite{Otto2001}, with the rigorous theory in \textcite{AmbrosioGigliSavare2008}. In machine learning, transport enters through distances between distributions \parencite{Cuturi2013}, as in the Wasserstein GANs of Arjovsky, Chintala, and Bottou (2017), through sampling, and through wide networks \parencite{ChizatBach2018,MeiMontanariNguyen2018}.

\printchapterbib
\end{refsection}
